\chapter{山东大学研究生学位论文撰写指南}

为提高学位工作水平和学位论文质量，保证学位论文在结构和格式上的规范与统一，依据《科学技术报告、学位论文和学术论文的编写格式》（GB/T 7713-1987）、《学位论文编写规则》（GB/T 7713.1-2006）、《信息与文献参考文献著录规则》（GB/T 7714-2015）、《科技文献的章节编号方法》（CY／T 35-2001）等相关国家标准，制定本指南。

\section{一、学位论文的基本要求}
学位论文一般应用中文撰写（外语专业除外），提倡并鼓励用中、外文撰写。各学科（专业学位类别）自行确定字数要求，博士学位论文字数一般不得低于3万字，摘要一般不超过2000字；硕士学位论文字数一般不得低于2万字，摘要一般不超过1000字。

\section{二、学位论文的结构要求}
博士、硕士学位论文一般应由以下几部分组成，依次为：
\begin{enumerate}
\item[(1)] 论文封面；
\item[(2)] 扉页；
\item[(3)] 原创性声明和关于论文使用授权的声明；
\item[(4)] 中文摘要；
\item[(5)] 外文摘要；
\item[(6)] 论文中、外文目录；
\item[(7)] 符号说明；
\item[(8)] 论文正文；
\item[(9)] 附录、附图表；
\item[(10)] 引文出处及参考文献；
\item[(11)] 致谢；
\item[(12)] 攻读学位期间发表的学术论文目录；
\item[(13)] 学位论文评阅及答辩情况；
\item[(14)] 外文论文（可选）。
\end{enumerate}

\subsection{（一）论文封面}
采用研究生院统一印制的封面。封面的论文题目需要中、外文标示。论文题目不得超过30个汉字。分类号须采用《中国图书资料分类法》进行标注。

\subsection{（二）扉页}
论文设扉页，其内容与封面相同，送交校学位办公室、图书馆和档案馆的论文，其扉页由本人用碳素钢笔或签字笔填写。

\subsection{（三）原创性声明和关于学位论文使用授权的说明}
论文作者和指导教师在向校学位办公室、图书馆、档案馆提交论文时必须在要求签名处签字。

\subsection{（四）中文摘要}
中文摘要应以最简洁的语言介绍论文的内容要点，并注意突出论文中的新论点、新见解或创造性的成果，在摘要后列出3～8个关键词，关键词之间用分号相隔。关键词应体现论文的主要内容，词组符合学术规范。

\subsection{（五）外文摘要}
外文摘要内容应与中文摘要基本一致，要语句通顺，语法正确，准确反映论文的内容，并在其后列出与中文相对应的外文关键词。

\subsection{（六）论文目录}
论文需要有中外文目录各一份。目录通常显示到三级，并注明页码。标题应简明扼要。

\subsection{（七）符号说明}
介绍论文中所用符号表示的意义。

\subsection{（八）论文正文}
正文是学位论文的主体和核心部分。正文部分由于涉及的学科、选题、研究方法、结果表达方式等有很大的差异，不能作统一的规定。但是，必须实事求是、客观真切、合乎逻辑、层次分明、简练可读。论文中的计量单位、制图、制表、公式、缩略词和符号等必须遵循国家规定的标准。其行文方式和文体的格局，研究生可根据自己研究课题的表达需要不同而变化，灵活掌握。章节标题应简明扼要，体现阐述内容的重点，无标点符号。

\subsection{（九）附录、附图表}
主要列入正文内过分冗长的公式推导，供查读方便所需的辅助性数学工具或表格；实验性图片；程序全文及说明等。

\subsection{（十）注释及参考文献}
注释是对论文正文中某一特定内容的进一步解释、补充说明或引文。一般在人文科学和社会科学中用的较多。注释体例应全篇统一。由于论文篇幅较长，建议注释采用脚注（每页重新编号，前后页不连续编号）。如确有需要，也可以采用尾注，将其置于于正文末尾。参考文献按文中使用的顺序列出，并注明文献的责任者、题名、刊物（出版社）名称、出版时间、页码等。

\subsection{（十一）致谢}
系对给予各类资助、指导和协助完成研究工作以及提供各种对论文工作有利条件的单位和个人表示的感谢。致谢应实事求是，切忌浮夸之词。

\subsection{（十二）攻读学位期间发表的学术论文目录}
按学术论文发表的时间顺序，列出本人在攻读学位期间发表或已录用的主要学术论文清单，包括顺序号、论文题名、刊物名称、卷册号及年月、起止页码、论文署名位次。

\subsection{（十三）学位论文评阅及答辩情况}
论文答辩通过后，送校学位办公室、图书馆和档案馆的论文需将学位论文评阅及答辩情况填入《学位论文评阅及答辩情况表》中。

\subsection{（十四）外文论文（可选）}
\subsubsection*{1．外文论文写作的形式}
可根据本学科的实际选择以下写作形式的其中一种。
\begin{enumerate}
\item[①] 与中文全文在内容和形式上完全一致的外文全文。
\item[②] 与学位论文相关的可以在外文期刊上发表（含已发表）的外文论文。
\end{enumerate}
\subsubsection*{2．外语写作的要求}
学位论文外语写作要语句通顺，语法正确，符合该种语言的写作规范，能准确反映作者的学术思想。论文内容用小四号字体小写字母打印。

\section{三、学位论文的格式要求}
\subsection{（一）所有英文字母、阿拉伯数字和半角标点符号默认用Times New Roman字体}
\subsection{（二）论文封面}
用小二号加粗黑体字打印封面的中文论文题目，用三号加粗打印封面外文论文题目，四号加粗黑体字打印脊背处论文题目和封面作者姓名、专业、指导教师、合作导师姓名和专业技术职务、论文完成时间、密级、学校代码、学号、分类号等内容。

\subsection{（三）中文摘要}
“摘要”标题字用小三号加粗黑体字打印，居中无缩进，段前24磅，段后18磅，单倍行距。

摘要内容用小四号宋体打印，两端对齐，首行缩进2字符，段前0行，段后0行，固定值20磅。

“关键词：”宋体加粗小四，其后关键词用全角分号分隔，中文字体为宋体，小四号，两端对齐，无缩进，段前0行，段后0行，固定值20磅。

\subsection{（四）外文摘要}
“ABSTRACT”标题字用加粗小三号字体打印，居中无缩进，段前24磅，段后18磅，单倍行距。

摘要内容用小四号字体打印，两端对齐，首行缩进2字符，段前0行，段后0行，固定值20磅。

“Key words:”加粗小四号，其后关键词用半角分号分隔，小四号，两端对齐，无缩进，段前0行，段后0行，固定值20磅。

\subsection{（五）论文目录}
中文的“目录”标题字用小三号加粗黑体字打印，居中无缩进，段前24磅，段后18磅，单倍行距。目录内容用小四号宋体打印，两端对齐，段前6磅，段后0磅，单倍行距。下级标题比上级左边多空两个汉字符宽度。外文目录的标题字用加粗小三号字体打印，目录内容用小四号字体打印。

\subsection{（六）论文正文}
宋体小四号，两端对齐，首行缩进2字符，段前0行，段后0行，固定值20磅（段落中有数学表达式时，可根据需要设置该段的行距）。引文内容可用楷体。

\subsection{（七）章、节标题}
理学、工学、医学类学位论文的章节编排建议遵循《CY／T 35-2001 科技文献的章节编号方法》，以阿拉伯数字编号，如1，1.2，1.2.1逐级递推。人文社科类学位论文建议采用第一章 第一节 一、(一)形式编排。英文撰写的论文建议采用Chapter1,1.2,1.2.1逐级递推。

章标题：黑体小三号加粗，居中无缩进，段前24磅，段后18磅，单倍行距。每一章另起一页。

二级标题：黑体四号加粗，两端对齐，无缩进，段前24磅，段后6磅，单倍行距。

三级标题：黑体小四号加粗，两端对齐，无缩进，段前12磅，段后6磅，单倍行距。

\subsection{（八）图}
图应具有“自明性”，即只看图、图题和图例，不阅读正文，就可理解图意。

图应分章编号，如图1-1。图宜有图题，图题即图的名称，置于图的编号之后。图的编号和图题应置于图下方。如有英文标题，以Fig. 1-1编号。图题用宋体五号，居中无缩进，段前6磅，段后6磅，单倍行距。图序与图名之间空一格，图名后不加标点。全文编号方式应统一。如果一个图由两个或两个以上分图组成时，各分图分别以(a)、(b)、(c)……作为图序，并须有分图名，分图名直接列在各自分图的正下方，主图名列在所有分图的下方正中。

可根据需要加图注（即图的注解和说明）。图注用阿拉伯数字按顺序编排，一般排在图题下面。

\subsection{（九）表}
表应具有“自明性”。表应分章编号，如表1-1；表宜有表题，表题即表的名称，置于表的编号之后。表的编号和表题应置于表上方。

表的编排一般采用国际通行的三线表。如某个表需要转页接排，在随后各页上应重复表的编号，如“续表2-1”，且续表均应重复表头。如有英文标题，以Tab. 1-1编号。全文编号方式应统一。

表题用宋体五号，西文Times New Roman，五号，居中无缩进，段前6磅，段后6磅，单倍行距。表序与表名之间空一格，表名后不加标点。

表中若有附注，用阿拉伯数字按顺序编排，附注写在表的下方。

\subsection{（十）公式}
论文公式应分章编号，并将编号置于括号内，如(1-1)，公式编号右端对齐，与公式间以...连接。包含公式的段落行间距可设置为最小值20磅，以保证公式完全显示。

\subsection{（十一）页眉页脚}
从摘要开始每页要有页眉，其上居中打印“山东大学博（硕）士学位论文”字样。页眉采用宋体五号字居中书写。页码从摘要开始，前置部分（摘要，ABSTRACT，目录等）用大写罗马数字（Ⅰ，Ⅱ，Ⅲ，……）编排，正文部分按阿拉伯数字（1，2，3，……）重新开始编号。页码位于页面底端，采用五号 Times New Roman居中书写。页码数字两侧不要加“-”等修饰线。

\subsection{（十二）参考文献}
参考文献标题：黑体小三号加粗，居中无缩进，段前24磅，段后18磅，单倍行距。另起一页。参考文献文本宋体小四号，两端对齐，段前0行，段后0行，固定值20磅。

\section{四、参考文献著录规则}
理工科参考文献著录项目和格式原则上应遵照《信息与文献参考文献著录规则》（GB/T 7714—2015）的规定执行。人文社科类学科可按通行惯例指定著录格式，如法学可参考北京大学出版社2020年5月出版的《法学引注手册》，但论文全篇必须保持统一。引证外文文献，原则上应以该文种通行的引证标注体例为准。《信息与文献参考文献著录规则》（GB/T 7714—2015）中几种主要类型的参考文献的著录项目与格式要求示例如下：

\subsection{（一）专著（图书）[M]：}
[序号]著者．题名：其他题名信息[M]．其他责任者．版本项．出版地：出版者，出版年：页码．

\begin{enumerate}
\item[1] 曹凌．中国佛教疑伪经综录[M]．上海：上海古籍出版社，2011：19.
\item[2] 钱学森．创建系统学[M]．太原：山西科学技术出版社，2001：序2-3．
\item[3] 冯友兰．冯友兰自选集[M]．2版．北京：北京大学出版社，2008：第1版自序．
\end{enumerate}

\subsection{（二）期刊论文[J]：}
[序号]作者．文献名[J]．期刊名，年，卷（期）：页码．

\begin{enumerate}
\item[1] 李炳穆．韩国图书馆法[J]．图书情报工作，2008，56(2)：6-12．
\item[2] 袁训来，陈哲，肖书海，等．蓝田生物群：一个认识多细胞生物起源和早期演化的新窗口[J]．科学通报，2012，55(34)：3219．
\end{enumerate}

\subsection{（三）学位论文[D]：}
[序号]作者．论文名[D]．学校所在城市：学校名，年份．

\begin{enumerate}
\item[1] 马欢．人类活动影响下海河流域典型区水循环变化分析[D]．北京：清华大学，2011．
\item[2] 赵睿智．通信工程学论文英译汉实践报告[D]．济南：山东大学，2017．
\end{enumerate}

\subsection{（四）报纸[N]：}
[序号]作者．题名[N]．报刊名，年-月-日(版数)．

\begin{enumerate}
\item[1] 丁文详．数字革命与竞争国际化[N]．中国青年报，2000-11-20(15)．
\item[2] 张田勤．罪犯DNA库与生命论理学计划[N]．大众科技报，2000-11-12(7)．
\end{enumerate}

\subsection{（五）论文集[C]：}
[序号]著者．论文集名[C]．出版地：出版者，出版年．

\begin{enumerate}
\item[1] 中国社会科学院台湾史研究中心．台湾光复六十五周年暨抗战史实学术研讨会论文集[C]．北京：九州出版社，2012．
\item[2] 牛志明，斯温兰德，雷光春．综合湿地管理国际研讨会论文集[C]．北京：海洋出版社，2012．
\end{enumerate}

\subsection{（六）标准文献[S]：}
[序号]标准制定者．标准名：标准号[S]．出版地：出版者，出版年：页码．

\begin{enumerate}
\item[1] 全国信息与文献标准化技术委员会．信息与文献 都柏林核心元数据元素集：GB/T 25100-2010[S]．北京：中国标准出版社，2010：2-3．
\end{enumerate}

\subsection{（七）专利[P]：}
[序号]专利所有者（申请者）．专利名：专利号[P]．公告日期．

\begin{enumerate}
\item[1] 张凯军．专利文献轨道火车及高速轨道火车紧急安全制动辅助装置：201220158825.2[P]．2012-04-05．
\end{enumerate}

\subsection{（八）档案、法律文件[A]：}
[序号]档案馆名．档案文献[A]．出版地：出版者，出版年．

\begin{enumerate}
\item[1] 中国第一历史档案馆，辽宁省档案馆．中国明朝档案总汇[A]．桂林：广西师范大学出版社，2001．
\end{enumerate}

\subsection{（九）报告[R]：}
[序号]主要责任者．题名：其他题名信息[R]．出版地：出版者，出版年份：页码．

\begin{enumerate}
\item[1] 中国互联网络信息中心．第29次中国互联网络发展现状统计报告[R]．北京：社会科学文献出版社，2012：84．
\end{enumerate}

\subsection{（十）专著中的析出文献}
析出文献指从整个信息资源中析出的具有独立篇名的文献，比如某论文集中的某一篇文章。其格式为：[序号]析出文献主要责任者．析出文献题名[文献类型标识]．专著主要责任者．专著题名：其他信息题名．版本项．出版地：出版者，出版年：析出文献的页码．（注意符号“//”表示“析出”）

\begin{enumerate}
\item[1] 周易外传：卷5[M]//王夫之．船山全书：第6册．长沙：岳麓书社，2011:1109．
\item[2] 马克思．政治经济学批判[M]//马克思，恩格斯．马克思恩格斯全集：第35卷．北京：人民出版社，2013：302．
\end{enumerate}

\subsection{（十一）网络资源：}
[序号]主要责任者．题名：其他题名信息[EB/OL]．（更新日期）[引用日期]．获取和访问路径．数字对象唯一标识符．

\begin{enumerate}
\item[1] 国务院学位委员会，教育部．博士硕士学位论文抽检办法[EB/OL]．(2014-01-29)．http://www.moe.gov.cn/srcsite/A22/s7065/201402/t20140212\_165556.html.
\end{enumerate}

\section{五、学位论文的打印与装订}
论文用A4标准纸输出，双面打印。页边距：上—2.8 cm，下—2.5 cm，左—2.5 cm，右—2.5 cm，装订线位置左侧，装订线0 cm。页眉边距1.5 cm，页脚边距1.75 cm。博士学位论文一式20份，硕士学位论文一式15份，装订成册，并按要求送交有关部门（送校图书馆和档案馆的论文需线装）。中、外文学位论文原则上一起装订，如篇幅过长可分别装订。除外语专业的学位论文外，其它学科的学位论文一律中文论文在前，外文论文在后。

\section{六、其他}
\begin{enumerate}
\item[(一)] 指南未尽之处，如有相关国家标准或通行惯例，遵循国家标准和通行惯例执行；如无标准或惯例，则根据具体情况调整，但应使论文整体风格一致。
\item[(二)] 为适应不同学科、不同类型学位论文特点，各培养单位可在参考本指南基础上，制订相应学位论文的具体写作规范，报学校批准备案后实施。如有必要，可在学位论文目录前增加说明页。
\item[(三)] 学校鼓励各专业学位类别积极探索符合本类别实际的学位论文规范。
\item[(四)] 本指南自发布之日起实施。原《山东大学学位论文规范（试行）》（山大研字〔2006〕53号）同时废止。
\item[(五)] 本指南由研究生院、党委研究生工作部负责解释。
\end{enumerate}